1: % Reference entries, D (AN-11).  Law 1: os/iword/core/edit.c
2: % iw_ed_backspace / iw_ed_delete (whole units; a hard return deleted joins
3: % two paragraphs and the joined paragraph is wrapped again), codes.c (the
4: % question for codes), os/iword/app/face_a.c edit_key (Backspace / Delete
5: % with any modifier act alone; harness/pubs/iword_probe.c lists them).
6: % GATES: test-pubs-iword test-iword-core test-iword-codes test-iword-face
7: \Gates{test-pubs-iword test-iword-core test-iword-codes test-iword-face}
8: \begin{Feature}{Delete}
9: You erase text with two keys: \K{Backspace} erases the character to the
10: left of the cursor, and \K{Delete} erases the character at the cursor.
11: Text after the cursor moves up to close the gap, and the paragraph is
12: wrapped again.
13: 
14: To delete text,
15: \begin{WPSteps}
16: \item Move the cursor just after the character you want to erase and press
17: \K{Backspace}
18: \Or
19: Move the cursor onto the character and press \K{Delete}.
20: \item Hold the key down to erase several characters in a row.
21: \end{WPSteps}
22: \Documents{wpkey}{Backspace}
23: \Documents{wpkey}{Delete}
24: \begin{WPNote}To delete more than a few characters, block them and press
25: Delete (see \InRef{Block, Delete}), or have Move delete the sentence,
26: paragraph or page (see \InRef{Move, Sentence/Paragraph/Page}).\end{WPNote}
27: 
28: \Notes
29: \Topic{Codes}
30: Backspace and Delete erase codes as they erase characters, one at a time.
31: Some codes are not erased at once on the editing screen: InitechWord asks
32: first (see \InRef{Codes}).
33: 
34: \Topic{Joining Paragraphs}
35: A hard return is erased like a character.  Erase the hard return at the end
36: of a paragraph -- with Delete at the end of its last line, or with
37: Backspace at the start of the next paragraph -- and the two paragraphs
38: become one.
39: 
40: \BeginReplay{ref-delete}
41: For example, type \TypeRun{First paragraph.}, press \PressKey{Enter}, type
42: \TypeRun{Second paragraph.}, press \PressKey{Home}, \PressKey{Left} and then
43: \PressKey{Backspace}:
44: \WordScreen[rows=1]{ref-delete}{Backspace at the start of a paragraph joins it to the one above}
45: \EndReplay
46: 
47: \Topic{Shift, Ctrl and Alt}
48: Backspace and Delete erase one character whatever Shift, Ctrl or Alt you
49: hold with them.
50: 
51: \Topic{The Start and End of the Document}
52: Backspace at the very start of the document and Delete at its very end do
53: nothing.
54: 
55: \SeeAlso{Block, Delete; Codes; Move, Sentence/Paragraph/Page}
56: \end{Feature}
57: 
58: \begin{Feature}{Directories}
59: InitechWord keeps your documents in one place: the top level of the
60: InitechOS data volume, the disk on which InitechWord itself is supplied.
61: InitechWord shows that place as \Scr{A:\wbs}, so that a document named
62: \Scr{MEMO} is \Scr{A:\wbs MEMO} on the status line and in its questions.
63: 
64: You do not type a drive or a directory when you save or retrieve: type the
65: name alone.  A name with a colon or a backslash is not a valid name (see
66: \InRef{Filenames}).
67: 
68: \Notes
69: \Topic{The Setup File}
70: InitechWord keeps its own settings in the same place, in the file
71: \Scr{IWORD.SET}, once you change one with Setup (see \InRef{Setup}).  Do
72: not save a document under that name: InitechWord would read it as its
73: settings, find it not in their form, and start with the standard ones; and
74: the next change you make in Setup would write over the document.
75: 
76: \SeeAlso{Filenames; Retrieve; Save; Setup}
77: \end{Feature}
